\begin{figure}[t]
    
    % \begin{subfigure}[b]{0.45\linewidth}
    %     \centering
    %     \includegraphics[width=\linewidth]{figs/SpatialArch-model.pdf} 
    % \end{subfigure}
  
    \begin{subfigure}[b]{\linewidth}
        \centering
        \includegraphics[width=0.5\linewidth]{figs/SpatialArch-folded.pdf}    
        \caption{Temporal architecture (i.e., overlay).}
        \label{subfig:temporal}
    \end{subfigure}
    % \vspace{1pt}
    
    \begin{subfigure}[b]{\linewidth}
        \centering
        \includegraphics[width=0.5\linewidth]{figs/SpatialArch-unfold-2.pdf}
        \caption{Partially unfolded spatial architecture with two PEs.}
        \label{subfig:partial}
    \end{subfigure}
    % \vspace{1pt}
    
    \begin{subfigure}[b]{\linewidth}
        \centering
        \includegraphics[width=0.5\linewidth]{figs/SpatialArch-unfold-4.pdf}    
        \caption{Fully unfolded spatial architecture with four PEs.}
        \label{subfig:spatial}
    \end{subfigure}
    % \vspace{1pt}
    
    \caption{Temporal and spatial architectures --- \textnormal{PE stands for processing engine; $f_1$-$f_4$ represent different operators in the model.}}
    \label{fig:temporal-spatial-arch}
\end{figure}
